% Measured diagnostic status, separate from proposed full method.
\subsection{条件反馈机制的诊断验证}
\label{sec:physics-guided-diagnostic-status}
为检验执行后条件误差能否改善参考构建，本轮冻结DiP生成器与CLoSD人体物理跟踪器，比较未经精化的生成参考、普通物理精化、等预算候选筛选、执行前缀反馈及有界条件纠偏，并分别去除、打乱反馈或更换前缀与筛选依据。每方法在48个原始MoCap来源、三个名义水平目标偏移和三个种子下进行432次独立回放；同一案例和种子内的原始目标固定，所有失败均保留。本轮是人体仿真诊断，不是目标人形机器人的部署实验，也不是Morph、PhysDiff或CLAIMS的完整复现。
普通精化、等预算随机筛选与候选条件纠偏的联合合格次数分别为0、0和0（各432次）。独立复核发现，参考和目标锚点在角色创建时构造，而实际角色随后重置，二者没有同步平移。种子变化因而同时改变生成噪声和目标相对实际角色的位置，名义目标偏移也混入了随机锚点差异。因此，初始化坐标测量校验没有通过；不能把这些数值视为受控偏移实验或纯随机种子重复，也不能据此证明条件纠偏有效或无效。
人体进入门槛未通过，因此按停止规则未开展G1迁移及下游策略训练。验证集还出现联合判据地板效应，主要受足部滑移代理与连续性限制；数值阈值在验证前固定，但没有按验证分布进行校准，这是与原计划的程序偏差。运行中还发生过实现核验与版本修订，早期输出的历史请求绑定不足。因此，本轮仅作为经过修订的诊断记录，不能作为确认性创新验证。
\begin{itemize}
\item 有效且有技术差异：本轮尚未建立此类结论。
\item 有效但属于已有机制：物理精化、执行反馈及闭环数据构建已有文献依据\citep{closd,parc,claims}；本轮不能仅凭回放完成率将其解释为联合质量改善。
\item 未获支持：本轮尚未取得可用于证明有界条件纠偏收益的有效实验依据；合成数据的下游训练收益尚未检验。
\end{itemize}
完整对比、消融、指标分解、尝试登记与版本记录保存在本项目的\texttt{literature\_review\_report.tex}研究记录及\href{https://github.com/Pendragon08/conditional-human-motion-generation}{公开实验仓库}中。后续须先独立核验任务窗口与验收判据的区分能力，再冻结新协议开展验证。

独立初始化复核仍未通过，不能据此展开机制有效性比较。独立复核不与旧轮合并，质量阈值未改变；首次修复在控制前的失败亦完整保留。
